% REVISED: canonical-target reframe 2026-06-04
% REVISED: hostile-review action plan 2026-06-04
% REVISED: JLEO positioning + humanization 2026-06-06
\section{Conclusion}
\label{sec:conclusion}

Cartel enforcement usually starts before an agency holds the
evidence a case would need. Procurement records then raise a problem
that precedes liability: how should scarce investigative capacity be
spent when the cheap administrative record is broad but legally thin,
and the richer bid-level record is costly to recover and evaluate?
We treat procurement-cartel screening as evidence allocation under
costly proof. The point is not to collapse the cheap award layer into
the costly bid layer but to sequence them, so administrative records
rank where investigation should begin and bid-level evidence judges
conduct inside the prioritized set.

%
The validation evidence maps that division of labor, and its sharp
limits are the substance of the finding. Persistent zero-win
participation orders adjudication-anchored loser-side exposure, but a
disciplined decomposition---opportunity-cell adjustment,
rolling-origin timing, leave-one-case-out, and environment-dominance
checks---shows that most of the cheap signal's apparent reach is
procurement opportunity and case concentration. Hold participation
volume and opportunity-set exposure fixed and the residual ordering
is marginal and fragile: near chance in matched cells, a small and
only borderline-significant nested increment, with matched-stratum
permutation and enrichment tests unable to reject the opportunity
null, and the performance riding on one adjudicated case and one
procurement environment. Prospective ranking fails outside incumbent
pools, since new entrants cannot be ordered by prior losing, leaving
the ranking retrospective and incumbent-bound. It does not behave like
a generic direct-defendant detector, exactly as its loser-side scope
implies. The general lesson is one of audit discipline: validating an
administrative screen against adjudicated cases without adjusting for
procurement opportunity systematically over-credits it. Separating
this fragile residual from exposure arithmetic and case
concentration---and saying plainly where each begins---is what the
method contributes.

The two layers also divide labor conditionally. On the same
adjudication-anchored target the award-layer score is comparable to a
transparent bid-moment benchmark, yet a combined model adds
information only under some validation folds and slips below the
award-layer ranking under case-grouped folds. The complementarity is
real but case-fragile, not a clean dominance. Joint scoring is a
full-observability upper bound; sequential gatekeeping is the
architecture that matters once an agency knows where the cheap layer
ranks and where it does not, and the cost--recall frontier---not any
single cutoff---is what it offers.

The limits are part of the design, not an apology for it. The cheap
screen creates forensic priority, not proof; the validation target is
adjudication-anchored exposure, not cartel membership; price evidence
is scope, not damages. The architecture travels only where award
records, repeated participation, comparable procurement categories,
and recoverable bid files coexist. Because a static cutoff can be
gamed, use should rest on refreshed thresholds, continuous ranks or
top-$k$ queues, bunching checks, and bid-layer follow-up.

The estimated ranking is not a portable cartel score; the transferable
object is the audit architecture---label construction, opportunity
adjustment, timing discipline, case-composition audit, bid-layer
comparison, and cost--recall accounting. We test that transferability
rather than assert it. Re-run unchanged on a second procurement
system---the federal ComprasNet platform, with partially overlapping
CADE anchors and no publicly recoverable bid layer---the same audit
returns the same verdict, deflating the cheap screen there as it does
on BEC. What ports across the two platforms is the discipline, not a
deployable ranking. Cheap suspicion should be neither ignored nor
mistaken for proof. Award-layer triage gives an agency a principled
way to order scarce attention before the richer evidentiary record is
assembled, and, just as usefully, a principled way to recognize when
an administrative signal is mostly exposure and concentration rather
than evidence. The statistic marks where legal attention should
start, not where legal responsibility ends; charting that reach, and
its edge, is what makes the cheap layer safe to work with.
